% Section 2: Reinforcement learning for forecasting
\subsection{Reinforcement learning for forecasting}
\label{sec:rl_for_forecasting}

\subsubsection{Reinforcement learning as an ensemble combiner}
\label{sec:rl_ensemble}

The simplest way to put reinforcement learning near forecasting is to treat the
dynamic combination of an ensemble of base forecasters as a sequential decision
problem. \citet{Fu2022} formalise this as a Markov decision process and call the
resulting agent RLMC. Given $N$ pretrained base forecasters
$\mathcal M = \{M_1, \dots, M_N\}$, the state at time $t$ is
$s_t \in \mathbb R^{T \times d_s}$, namely the recent $T$-step window of the time
series together with a log of recent per-model performance. The action
$a_t = (w^1_t, \dots, w^N_t) \in \Delta^{N-1}$ is a vector of non-negative weights
on the simplex, and the realised forecast at time $t$ is the convex combination
$\hat y_t = \sum_{i=1}^N w^i_t \hat y^i_t$ of the base predictions. The transition
satisfies
\begin{equation}
P(s_{t+1} \mid s_t, a_t) = P(s_{t+1} \mid s_t),
\label{eq:fu2022_transition}
\end{equation}
that is, the agent's weighting choice does not move the underlying time series.
The reward is a mixture of a sMAPE-based and a rank-based normalised score,
$r_t = \alpha R^{\mathrm{sMAPE}}_t + (1-\alpha) R^{\mathrm{rank}}_t$, where each
term is renormalised to $[-1, 1]$ by quantile transformation against the
in-sample distribution of base-model errors at time $t$.

\citet{Fu2022} draw two methodologically useful observations from this setup. The
first concerns exploration. Because the action does not enter the state-transition
kernel by~\eqref{eq:fu2022_transition}, the conditional next-state distribution can
be sampled from the empirical state-trajectory of the training data with any
action superimposed. The combination problem is therefore in effect a model-based
RL problem in which the transition kernel and the reward are both known up to the
empirical sample, and an unbounded supply of transitions
$(s_t, a_t, r(s_t, a_t), s_{t+1})$ can be generated without further interaction. The
second observation concerns the structure of the optimal policy. If only one base
model is best at each state, the optimal policy is nearly deterministic, namely
$\pi^\star(s) \approx e_{i^\star(s)}$, the standard basis vector at the best model
index. A policy-gradient method then drives $\log \pi_\phi(a \mid s) \to -\infty$
on most of the action space, and the estimator
$\hat g = \mathbb E_t[\nabla_\phi \log \pi_\phi(a_t \mid s_t) \hat A_t]$ becomes
numerically unstable. Figure 2 of \citet{Fu2022} documents the divergence
empirically on the ETT dataset.

These observations motivate the choice of Deep Deterministic Policy Gradient
\citep{LillicrapDDPG2016} as the optimisation algorithm. The actor
$a = \pi_\phi(s)$ outputs a continuous deterministic action, and the critic
$Q_\theta(s, a)$ approximates the optimal action-value function under the
mean-squared Bellman loss
$\mathcal L_Q(\theta) = \mathbb E_{(s,a,r,s')}[ (Q_\theta(s,a) - r - \gamma
Q_{\theta'}(s', \pi_\phi(s')))^2 ]$, with the policy updated by deterministic
policy gradient
$\nabla_\phi J = \mathbb E_s[ \nabla_a Q_\theta(s, a)\rvert_{a = \pi_\phi(s)}
\nabla_\phi \pi_\phi(s) ]$. The deterministic actor sidesteps the log-probability
explosion, and the off-policy structure exploits the model-based generation route
in the first observation. Three exploration tricks complete the algorithm. The
actor is pretrained on a multi-class classification surrogate whose labels are
the indices of the per-state best base model. An $\epsilon$-greedy rule with a
convex-combination inductive bias replaces the standard Gaussian perturbation,
namely $a = \pi_\phi(s)$ with probability $1 - \epsilon$ and $a = e_i$ for a
uniformly drawn $i$ with probability $\epsilon$. A second replay buffer holds
hard transitions whose realised reward falls below a threshold $\hat r_t$, and
mini-batches are drawn from both buffers so that the easy distribution does not
absorb the agent toward a degenerate copy-the-best-model policy.

The empirical evidence in \citet{Fu2022} covers four datasets. On the ETT
electricity-transformer-temperature series, RLMC reports a sMAPE of $72.11$ against
$73.69$ for the single best base model and $79.83$ for the uniform average. On
the M4 daily panel, RLMC reports a sMAPE of $2.31$ against $2.39$ for the single
best and $3.93$ for uniform. On the Climate series the MAE falls to $1.47$ against
$1.67$ for single best, and on the GEFCOM electricity-load benchmark RLMC ties
with FFORMA at MAE $88$ to $87$. The gains are real but modest, in the range of
one to three per cent of sMAPE or MAE relative to the single best base model,
and are sometimes overtaken on individual datasets by the gradient-boosted
FFORMA baseline. The ablation in Figure 7 of \citet{Fu2022} suggests that the
exploration strategy carries most of the improvement, with the pretraining and
the hard-sample buffer contributing marginal additional gains.

The methodological depth of this strand is modest, and we mention it primarily to
demarcate the scope of the chapter. The model combination problem in
\citet{Fu2022} can be recast as an online learning problem on the simplex with a
non-stationary loss sequence. The regret guarantees one would obtain from a Hedge
or an EXP3 algorithm against the best fixed combination in hindsight are
$O(\sqrt{T \log N})$ and $O(\sqrt{N T \log N})$ respectively, and there is no
indication that DDPG-based weight learning improves on these rates. The genuine
contribution of \citet{Fu2022} is the demonstration that off-policy actor-critic
machinery can be applied to model-combination problems with imbalanced
best-model distributions, not a regret bound. The deeper RL-for-forecasting
strands, namely Brier-loss reward training, approximate dynamic programming on a
sequential-estimation MDP, and decision-focused learning, sit on richer
methodological foundations and occupy the remainder of this section.

\subsubsection{Reinforcement learning with decision-relevant rewards}
\label{sec:rl_brier}

A second strand of RL-for-forecasting uses on-policy reinforcement learning to fine
tune a probabilistic forecaster against a proper scoring rule that doubles as the
downstream decision loss. The cleanest production-scale instance to date is
\citet{Turtel2025}, who adapt the family of reinforcement-learning-with-verifiable-rewards
algorithms, which had previously been confined to deterministic mathematical
verification, to the stochastic and lagged-supervision setting of real-world
forecasting. The base model is the 14B-parameter DeepSeek-R1-Distill-Qwen reasoning
model, the training corpus consists of $10{,}000$ historical Polymarket yes/no
questions augmented with $100{,}000$ synthetic questions generated by the Lightning
Rod Foresight framework, and the evaluation is a held-out set of $1{,}265$ test
questions with strict precautions against temporal leakage.

The forecasting episode is a single decision. The agent receives a textual prompt
that includes the question and a set of news headlines dated strictly before the
prediction date, and emits a single probability $\hat p \in [0, 1]$ for the event.
After the resolution date, the binary outcome $y \in \{0, 1\}$ is revealed and the
episode terminates with reward $R = -(\hat p - y)^2$. The reward is the negative
Brier score of \citet{Brier1950}, a strictly proper scoring rule in the sense of
\citet{GneitingRaftery2007}, so the population-optimal policy is the true
conditional probability $\hat p = \mathbb P(y = 1 \mid \text{prompt})$. The training
schedule is strictly online and single-pass, with questions encountered in
chronological order and outcomes revealed at the resolution date. Multiple epochs
are not used because re-exposure to a resolved question lets the model recover the
future from training-time information.

The algorithmic contribution is to identify why standard GRPO destabilises in this
regime and how to fix it. Group Relative Policy Optimisation \citep{Shao2024deepseek}
draws $G$ rollouts $\{o_1, \dots, o_G\}$ from the old policy and forms the advantage
$\hat A_i = (r_i - \mu) / \sigma$, where $\mu$ and $\sigma$ are the mean and
standard deviation of the rewards across the group. The per-question $\sigma$
normalisation is helpful when rewards are scale-heterogeneous across questions, but
\citet{Turtel2025} argue that it is harmful when the reward is a Brier score because
the asymmetry between modest correct-and-overconfident gains and rare large
miscalibration penalties is precisely the signal the model needs to learn proper
calibration. Their Modified-GRPO drops the $\sigma$ division and uses the
baseline-subtracted advantage $\hat A_i = r_i - \mu$, paralleling the Dr.\,GRPO
modification of \citet{Liu2025DrGRPO}. ReMax \citep{LiReMax2023} replaces the group
mean with a learned baseline $b_i$ and uses $\hat A_i = r_i - b_i$, which removes
the variance normalisation while preserving an unbiased advantage estimate. A
malformed model output that fails the regex extraction is assigned the maximum
Brier loss of $1.0$ during training (strict Brier) and a constant $0.25$ during
evaluation (soft Brier), with the latter chosen to preserve gradient signal without
inducing training collapse.

The failure modes that motivated these modifications are documented qualitatively
and quantitatively. Standard GRPO without guardrails drives the model toward
extreme overconfidence, with $39.3\%$ of probabilities falling into the $0$-$10\%$
or $90$-$100\%$ buckets and many landing on exact zero or one. Modified-GRPO
without guardrails reduces the extreme fraction to $7.9\%$ but introduces a
different pathology, namely code-switching between Chinese and English mid-sequence
and ungrammatical partial outputs. Adding the strict-Brier reward, a token-length
cap, an early-stopping gibberish filter, and the ReMax baseline brings the extreme
fraction down to $13.1\%$ while restoring grammatical coherence. The optimisation
scaffold across algorithms is AdamW with the standard reasoning-LLM hyperparameters
and bfloat16 precision, with global gradient-norm clipping at $1.0$ and an entropy
bonus coefficient of $0.001$. Total compute is roughly three days on eight NVIDIA
H100 GPUs per training run.

The headline empirical results vindicate the design. On the $1{,}265$-question hold
out, the seven-run ReMax ensemble trained on the $110{,}000$-question combined
corpus attains a soft-Brier score of $0.190$ with $95\%$ confidence interval
$[0.178, 0.203]$ and an expected calibration error of $0.062$, against $0.202$ and
$0.093$ for OpenAI o1, $0.215$ and $0.089$ for the base DeepSeek-R1 14B model, and
$0.151$ and $0.043$ for the contemporaneous Polymarket market price. The Brier
improvement over o1 is $-0.011$ with two-sided $p < 0.05$, and the calibration
improvement is large in absolute terms but does not reach statistical significance
against Polymarket itself. None of the trained models match the market on Brier,
which is consistent with the market aggregating information across many traders,
but the trained forecasters retain incremental information beyond the price.

The decision-side evaluation establishes the practical relevance. \citet{Turtel2025}
construct a hypothetical trading rule that places a one-share bet on every test
question and measure cumulative profit on Polymarket. The ReMax ensemble earns
$\$52$ in simulated profit on $\$433$ of committed capital, corresponding to
roughly a $12\%$ return on investment. The Modified-GRPO model earns $\$54$ on
$\$428$, and OpenAI o1 earns $\$39$. The profit concentrates on questions where
the market was most uncertain. On the subset with market probability between $40\%$
and $60\%$, the ReMax ensemble's bets win $11.8$ percentage points more often than
the market price predicts, with $95\%$ confidence interval $[8.0, 15.7]$ and a
return on investment of roughly $20\%$. The pattern is the empirical realisation of
the Granger-Pesaran argument from
Section~\ref{sec:forecasting_rl_intro}. A forecaster trained to minimise a proper
scoring rule, with the reward instrument set to the same Brier loss that the
evaluator computes, recovers calibration even when the prediction-accuracy
performance does not surpass an aggregating market, and the calibration translates
into realised profit on the threshold-crossing decisions where the market is
indifferent.

The methodological lesson generalises beyond LLMs and prediction markets. Treating
a strictly proper scoring rule as the reward signal of a sequential RL agent
converts the calibration question into a policy-optimisation question, which is
amenable to standard variance-reduction techniques. The two failure modes
identified by \citet{Turtel2025}, namely extreme overconfidence under raw GRPO and
gibberish under raw Modified-GRPO, are not specific to LLMs but to the interaction
of variance-normalised policy gradients with bounded reward distributions. A
forecasting workflow that emits continuous quantile predictions and trains under
the continuous ranked probability score of \citet{GneitingRaftery2007}, or that
emits a categorical distribution and trains under the log score, can be expected
to encounter the same gradient instabilities and to benefit from the same
baseline-subtraction modifications. The connection between Brier-loss RL training
and the SPO programme of Section~\ref{sec:dfl} is direct. Both are special cases
of the differentiable-Wald-decision-theory recipe, with the difference being that
Brier-loss RL trains under the scoring rule itself rather than under the realised
decision regret of a specific downstream optimisation.

\subsubsection{Approximate dynamic programming for forecasting}
\label{sec:adp_forecasting}

Sequential estimation has a longer pedigree in the dynamic programming tradition
than in the deep learning tradition. \citet{BertsekasRollout2022} provides the
unifying frame and the vocabulary that the remainder of this subsection uses. The
sequential estimation problem is cast as a Markov decision process whose state at
time $k$ is the posterior $b_k$ over an unknown $m$-dimensional parameter
$\theta = (\theta_1, \dots, \theta_m)$, whose action is the index $u_{k+1} \in
\{1, \dots, m\}$ of the next noisy measurement $z_{u_{k+1}} = \theta_{u_{k+1}} +
w_{u_{k+1}}$, and whose transition follows the Bayes recursion
$b_{k+1} = B_k(b_k, u_{k+1}, z_{u_{k+1}})$. Per-observation costs $c(u_k)$ and a
terminal cost $G(b_N)$ that depends on the posterior mean and covariance complete
the cost structure. The exact dynamic programming algorithm
\begin{equation}
J^\star_k(b_k) = \min_{u_{k+1} \in \{1, \dots, m\}}
  \Bigl\{ c(u_{k+1}) + \mathbb E_{z_{u_{k+1}}}\bigl[
    J^\star_{k+1}\bigl(B_k(b_k, u_{k+1}, z_{u_{k+1}})\bigr)
  \bigr] \Bigr\}
\label{eq:bertsekas_dp}
\end{equation}
is infeasible to execute in general because the space of posteriors is
infinite-dimensional, and it is impractical even in the Gaussian special case
where the posterior reduces to a mean and a covariance.

The remedy is approximation in value space. The myopic acquisition-function policy
familiar from Bayesian optimisation, $\hat u_{k+1} = \arg\max_u A_k(b_k, u)$, is
one such approximation, and the rollout algorithm of \citet{BertsekasRollout2022}
is another. Given a base policy $\mu_{k+1}(\cdot)$ whose cost function
$\widetilde J_{k+1}$ can be evaluated by Monte Carlo simulation, the rollout
controller selects
\begin{equation}
\widetilde u_{k+1} = \arg\min_{u_{k+1}} \Bigl\{
  c(u_{k+1}) + \mathbb E_{z_{u_{k+1}}}\bigl[
    \widetilde J_{k+1}\bigl(B_k(b_k, u_{k+1}, z_{u_{k+1}})\bigr)
  \bigr] \Bigr\}.
\label{eq:rollout_update}
\end{equation}
The cost-improvement guarantee, namely that the rollout policy weakly dominates
the base policy under mild contraction conditions, is the formal underpinning of
the entire ADP-for-forecasting programme. Certainty equivalence further reduces
the Monte Carlo burden by replacing the stochastic continuation cost with its
expectation conditional on the realised first step, at the price of a controlled
loss of optimality. The framework subsumes Bayesian optimisation, sequential
decoding under a deterministic observation channel (with Wordle and Mastermind as
illustrative cases), and adaptive control of stochastic systems with unknown
parameters, and it returns the same Bellman-equation language that the remainder
of this chapter exploits.

\citet{Fernandes2024} applies the neuro-dynamic programming variant of this
machinery to a quarterly gross domestic product nowcasting problem. The Bellman
equation is
\begin{equation}
J(i) = \min_u \mathbb E\bigl[ g(i, u, j) + \alpha\, J(j) \bigr],
\label{eq:fernandes_bellman}
\end{equation}
where $i$ is the current state of the economy, summarised by the regularised
levels of industrial production, construction output, retail trade turnover,
goods exports, goods imports, and the European Commission's Economic Sentiment
Indicator. The control $u$ is the quarter-over-quarter change of output, the cost
$g(i, u, j)$ is the quadratic adjustment cost between successive output levels in
the spirit of Lucas-Prescott (1971), and $\alpha \in (0, 1)$ is the discount
factor. The reward is therefore not the cross-sectional forecast error of OLS,
but the temporal difference between successive predictions of the level, a
reflection of how a central bank or finance ministry actually uses a nowcast.
\citet{Fernandes2024} approximates $J^\star$ by a single-hidden-layer feedforward
network with ReLU activations and tensor-product features and trains the network
by Sutton's TD(0) algorithm with the update
$\delta \leftarrow \widetilde J(i, r) - \alpha \widetilde J(j, r) - g(i, u, j)$
and gradient-style weight steps. Training uses a panel of $26$ European Union
countries, $2{,}470$ quarterly transitions across the period 2000Q1 to 2023Q4,
with Portugal held out for evaluation over 2015Q1 to 2023Q4.

The empirical findings of \citet{Fernandes2024} are striking for an economics
audience. On the regularised Portuguese GDP level, the OLS benchmark fitted on
Portugal-only data 2000-2014 attains a mean absolute error of $6.61\%$ of the
GDP range and a root mean squared error of $8.32\%$, respectively. The TD-trained
neural network on the EU-ex-Portugal panel attains an MAE of $4.53\%$ and an RMSE
of $5.97\%$, corresponding to a $32\%$ reduction in MAE and a $28\%$ reduction in
RMSE relative to OLS. A linear architecture trained with TD on the same EU panel
fares worst at $10.70\%$ MAE and $12.81\%$ RMSE, which establishes that the gain
is attributable to the nonlinear approximation and not to the panel-data
augmentation alone. The COVID-19 quarter is the canonical hard case, and the
cumulative absolute out-of-sample error figure in \citet{Fernandes2024} shows
that the TD-trained nonlinear network tracks the spring-2020 collapse without the
large prediction error that OLS suffers. The most surprising methodological
finding is that training on foreign data transfers cleanly to the held-out
country, which alleviates the small-$T$ problem that plagues macroeconomic
forecasting and supplies the most concrete macro-domain validation of NDP/RL on
record.

\citet{Capitaine2025} extend the dynamic programming side of the literature into
the online and non-stationary regime, providing the first provable guarantees for
online decision-focused learning. Their setting is a bilevel optimisation problem
over $T$ rounds. At round $t$, nature picks a covariate $X_t$ and a cost function
$\bar g_t(\cdot)$, the decision-maker emits a prediction $g(\theta_t, X_t)$ and
takes the action $w^\star_t(\theta_t) = \arg\min_{w \in \mathcal W}
\langle g(\theta_t, X_t), w \rangle$ over a polytope $\mathcal W$, observes the
realised cost $\bar g_t(X_t)$, and updates $\theta_{t+1}$. The induced
single-round objective $f_t(\theta) = \langle \bar g_t(X_t), w^\star_t(\theta)
\rangle$ is not differentiable because $\theta \mapsto w^\star_t(\theta)$ ranges
in the finite vertex set of $\mathcal W$. \citet{Capitaine2025} regularise the
inner problem by replacing $w^\star_t(\theta)$ with
$\widetilde w_t(\theta) = \arg\min_{w \in \mathcal W}
\{ \langle g(\theta, X_t), w \rangle + \alpha_t R(w) \}$
under a strongly convex regulariser $R$, namely the log-barrier on a general
polytope or the negative entropy on the simplex. The smoothed surrogate
$\widetilde f_t$ is Lipschitz but generically non-convex, which forces the use of
an $\xi$-approximate offline optimisation oracle in place of exact minimisation.

\citet{Capitaine2025} present two algorithms. Decision-Focused Follow-the-Perturbed
-Leader (DF-FTPL) selects $\theta_{t+1}$ by oracle minimisation of the perturbed
cumulative surrogate $\sum_{s \le t} \widetilde f_s - \langle \sigma_t, \cdot
\rangle$ with $\sigma_t \sim \mathrm{Exp}(\eta)^m$, and achieves a static regret
bound of order $\widetilde O\bigl( m^{3/4} n^{1/2} T^{-1/4} \bigr)$ against the
best fixed parameter in hindsight, with the dimension of the decision space
$\mathcal W$ entering only through a $\ln \ln d$ factor. Decision-Focused Online
Gradient Descent (DF-OGD) updates $\theta_{t+1}$ by a single gradient step at a
uniformly drawn convex combination of the current iterate and the per-round oracle
minimiser, and achieves a dynamic regret bound of order
$\widetilde O\bigl( (1 + P_T)^{1/4} T^{-1/4} n^{-1/2} \bigr)$, where
$P_T = \sum_{t < T} \lVert \vartheta_{t+1} - \vartheta_t \rVert$ is the path
variation of the approximate minimisers. Both rates are sublinear, both are
independent of the parameter dimension up to logarithmic factors, and both
explicitly require the margin condition of \citet{Mammen1999} and
\citet{Tsybakov2004}, namely that the optimal vertex of $\mathcal W$ is
identifiable in the sense that the gap between the optimal vertex and the second
best satisfies a polynomial probability bound. The knapsack-style experiment of
\citet{Capitaine2025} on a synthetic non-stationary cost
$\bar g_t(X_t) = \sin^4((2 X_t^\top \theta_t^\star)^{-1})$ with a misspecified
linear predictor shows DF-FTPL and DF-OGD outperforming prediction-focused OGD
and online SPO on cumulative cost, at the expected price of a higher mean squared
prediction error.

\subsubsection{Decision-focused learning}
\label{sec:dfl}

The deepest sub-stream of RL-for-forecasting is decision-focused learning, which
trains the predictor under a loss that reflects the downstream decision regret
rather than the cross-sectional prediction error. The programme was formalised by
\citet{ElmachtoubGrigas2022} under the name Smart Predict-then-Optimize, and the
setup it introduces governs the remainder of this subsection. A contextual
stochastic optimisation problem has a known feasible region $\mathcal Z \subseteq
\mathbb R^d$, a feature vector $x \in \mathcal X$, and a cost vector $c \in
\mathbb R^d$ whose conditional distribution given $x$ is $\mathcal D_x$. The
decision maker observes $x$, predicts $\hat c$ from a hypothesis class
$\mathcal H$, and then solves the deterministic problem
$\min_{z \in \mathcal Z} \hat c^\top z$ at decision time. The realised quality of
the prediction is measured by the SPO loss, namely
\begin{equation}
\ell_{\mathrm{SPO}}(\hat c, c) = \max_{z \in W^\star(\hat c)} c^\top z - z^\star(c),
\label{eq:spo_loss}
\end{equation}
where $W^\star(\hat c) = \arg\min_{z \in \mathcal Z} \hat c^\top z$ and
$z^\star(c) = \min_{z \in \mathcal Z} c^\top z$ is the optimal value under the
true cost. The pointwise maximum over $W^\star(\hat c)$ disambiguates the loss
when the inner optimisation has multiple optimisers. Proposition~1 of
\citet{ElmachtoubGrigas2022} establishes that on the simplest non-trivial
feasible region $\mathcal Z = [-1/2, 1/2]$ with $c \in \{-1, +1\}$, the SPO loss
specialises to the $0$-$1$ classification loss, which means SPO sits one structural
generalisation above the canonical loss of supervised binary classification.

The SPO loss is discontinuous and non-convex in $\hat c$, so empirical risk
minimisation is intractable. \citet{ElmachtoubGrigas2022} derive a convex
surrogate via Lagrangian duality and a first-order approximation, namely the
SPO$^+$ loss
\begin{equation}
\ell_{\mathrm{SPO}^+}(\hat c, c)
= \max_{z \in \mathcal Z}
  \bigl\{ c^\top z - 2 \hat c^\top z \bigr\}
  + 2 \hat c^\top z^\star(c) - z^\star(c),
\label{eq:spo_plus_loss}
\end{equation}
which is a convex upper bound on the SPO loss and admits the explicit subgradient
$2( z^\star(c) - z^\star(2 \hat c - c) )$ at $\hat c$. On the binary-classification
specialisation above, SPO$^+$ equals the hinge loss evaluated at twice the
prediction, which is the regulariser of choice in the support vector machine. The
SPO$^+$ ERM problem on a linear hypothesis class can be reformulated as a
linear or conic programme and solved by stochastic subgradient methods at scale.

The central theoretical claim is Fisher consistency, which justifies SPO$^+$ as
a population-level proxy for SPO. Let
$R_{\mathrm{SPO}}(\hat c) = \mathbb E_c [\ell_{\mathrm{SPO}}(\hat c, c)]$ and
$R_{\mathrm{SPO}^+}(\hat c) = \mathbb E_c [\ell_{\mathrm{SPO}^+}(\hat c, c)]$
denote the Bayes risks of the two losses.

\begin{theorem}[\citealp{ElmachtoubGrigas2022}, Theorem 1]
\label{thm:spo_fisher}
Suppose, conditional on $x$, that $W^\star(\mathbb E[c \mid x])$ is a singleton
almost surely, the distribution of $c \mid x$ is continuous on $\mathbb R^d$ and
centrally symmetric about its mean, and the interior of $\mathcal Z$ is nonempty.
Then any minimiser of $R_{\mathrm{SPO}^+}$ over the unrestricted hypothesis class
equals $\mathbb E[c \mid x]$ almost surely and also minimises $R_{\mathrm{SPO}}$.
Equivalently, the SPO$^+$ loss is Fisher consistent with the SPO loss.
\end{theorem}

The central symmetry assumption is non-trivial and is exactly where the SPO$^+$
recipe trades off generality for tractability. The least-squares loss is Fisher
consistent under the same assumption set, so SPO$^+$ and least squares share the
asymptotic-optimality property under well-specification. The case for SPO$^+$
over least squares is therefore an argument about finite-sample behaviour and
about robustness to misspecification, not about asymptotic minimisation.

\citet{LiuGrigas2021} convert Fisher consistency into a finite-sample statement.
Their main technical object is the calibration function
$\delta_{\mathrm{SPO}^+}(\varepsilon)
= \inf \bigl\{ R_{\mathrm{SPO}^+}(\hat c) - R_{\mathrm{SPO}^+}^\star
  : R_{\mathrm{SPO}}(\hat c) - R_{\mathrm{SPO}}^\star \ge \varepsilon \bigr\}$,
which measures the minimum surrogate-risk gap implied by a given excess SPO risk.
They prove that for a convex polyhedral feasible region $\mathcal Z$, the
calibration function satisfies $\delta_{\mathrm{SPO}^+}(\varepsilon) \ge C
\varepsilon^2$ for a constant $C$ depending on the geometry of $\mathcal Z$ and on
moments of the noise. Combined with a standard Rademacher complexity bound on the
SPO$^+$ excess surrogate risk, this yields an excess SPO risk of order $n^{-1/4}$
for empirical risk minimisation over a Rademacher-bounded linear hypothesis class.
For strongly convex level sets, namely entropy-constrained or quadratic-constrained
feasible regions, the calibration function improves to $\delta_{\mathrm{SPO}^+}
(\varepsilon) \ge C \varepsilon$, and the excess SPO risk decays at the parametric
rate $n^{-1/2}$, matching that of squared-error regression. The empirical study
of \citet{LiuGrigas2021} on entropy-constrained portfolio allocation confirms the
faster rate.

\citet{DontiAmosKolter2017} operationalise the decision-focused programme for
neural networks. Given a parametric forecaster $p(y \mid x; \theta)$ and a
downstream stochastic programme $z^\star(x; \theta) = \arg\min_z
\mathbb E_{y \sim p(y \mid x; \theta)}[ f(x, y, z) ]$ subject to constraints, they
construct the task loss $L(\theta) = \mathbb E_{(x, y)}[ f(x, y, z^\star(x;
\theta)) ]$ and minimise it by stochastic gradient descent. The technical core is
the differentiation of the optimal-solution map $z^\star(x; \theta)$ through the
Karush-Kuhn-Tucker conditions of the inner programme. Differentiating the KKT
equations and applying the implicit function theorem yields a linear system whose
solution returns the Jacobian
$\partial z^\star / \partial \theta$, which is solved efficiently with the
OptNet-style QP-derivative routine. They validate the method on three economic
problems. On a synthetic conditional newsvendor with quadratic over- and
under-order costs, the task-based linear model dominates both maximum-likelihood
plug-in and pure policy optimisation when the underlying conditional distribution
is misspecified. On a 24-hour generator-scheduling task fitted to seven years of
PJM electrical-load data, the task-based neural network reduces realised
scheduling cost by $38.6\%$ relative to a maximum-likelihood-trained Gaussian
load model used in the same optimisation, and by $8.6\%$ relative to a
cost-weighted-RMSE variant. On a six-year grid-scale battery arbitrage task, the
task-based net improves on the maximum-likelihood net by between $2\%$ and $102\%$
across the relevant hyperparameter range, with the larger gains occurring at low
to moderate flexibility weights. The implicit-differentiation routine introduced
here is the methodological bridge that lets the rest of the DFL literature scale
from convex linear programmes to deep neural networks.

The DFL story is not uniformly positive, and the most important counterpoint is
\citet{HuKallus2020}. They study the contextual linear optimisation problem
$\pi^\star(x) = \arg\min_{z \in \mathcal Z} \mathbb E[Y \mid X = x]^\top z$ under
a polytopal feasible region. The induced empirical risk minimisation method of
\citet{ElmachtoubGrigas2022}, which integrates prediction and decision into one
loss, is contrasted with the naive estimate-then-optimise plug-in that fits
$\hat f(x) = \mathbb E[Y \mid X = x]$ by ordinary least squares and then plugs
the estimate into the deterministic problem. Under a margin condition that
bounds the probability that the optimal vertex of $\mathcal Z$ is within
$\varepsilon$ of the second-best vertex, namely
$\mathbb P[ \min_{i \ne i^\star} (u_i - u_{i^\star}) \le \varepsilon ] \le C
\varepsilon^\beta$ for some $\beta \in (0, 1]$, plug-in ETO achieves an average
regret rate of order $n^{-1}$ when $\beta = 1$, while the integrated method
achieves only $n^{-2/3}$. The plug-in exploits the curvature of the conditional
mean function via off-the-shelf nonparametric regression in a way that the
integrated objective discards. \citet{HuKallus2020} read this as a corrective
rather than a refutation, namely that the case for DFL must be argued against
the specific economic setting and the specific hypothesis class, not in
generality.

\citet{HuangGupta2024} address the gap that \citet{HuKallus2020} reveal, by
constructing decision-focused surrogates whose guarantees hold under
misspecification. Their Perturbation Gradient losses approximate the directional
derivative of the plug-in decision loss using zeroth-order finite-difference
estimation along randomly perturbed prediction directions. The PG family is
Lipschitz continuous, expressible as a difference of concave functions, and
optimisable by off-the-shelf stochastic gradient descent given only black-box
access to the inner optimiser. The central theoretical result states that under a
sub-Gaussian noise model and a Rademacher complexity bound $R_n$ on the
hypothesis class, the population minimiser of the PG loss attains an excess
regret of order $\widetilde O\bigl( R_n + n^{-1/2} \bigr)$ against the best
in-class policy, which for linear hypothesis classes simplifies to
$\widetilde O\bigl( (dp/n)^{1/4} \bigr)$. The bound is best-in-class rather than
full-information optimal, but it does not require realisability of the true
conditional mean by the predictor, which is the operative condition of
\citet{ElmachtoubGrigas2022}'s Fisher-consistency theorem. The empirical study
of \citet{HuangGupta2024} on a vehicle-routing and a portfolio problem confirms
that PG losses dominate SPO$^+$ under misspecification by a substantial
margin and that they match SPO$^+$ under well-specification, which positions PG
as the current methodological frontier for decision-focused learning.

\citet{Mandi2024} provide the consolidating survey and the canonical empirical
benchmark. They organise the field into gradient-based methods, including
surrogate-loss schemes such as SPO$^+$, smoothing-based schemes such as the
Fenchel-Young loss of Berthet et al.\ (2020), and implicit-differentiation schemes
such as OptNet and the QPTH solver, and gradient-free methods including
restricted-policy parametrisations and Bayesian-optimisation-based meta-learning.
Their benchmark covers eleven DFL techniques on seven canonical problems, namely
shortest path, weighted bipartite matching, energy-cost-aware scheduling,
portfolio optimisation, knapsack, weighted travelling salesman, and a warcraft
combinatorial map task. The headline finding is that no single DFL technique
dominates on all seven problems, that the empirical advantage of DFL over
prediction-focused ERM is largest precisely under misspecification, and that the
relative ranking among DFL techniques is sensitive to the structural properties
of the inner programme. The composite picture across
\citet{ElmachtoubGrigas2022}, \citet{LiuGrigas2021},
\citet{DontiAmosKolter2017}, \citet{HuKallus2020}, \citet{HuangGupta2024}, and
\citet{Mandi2024} is therefore that DFL is the right asymptotic objective under
misspecification, that the rate is at best $n^{-1/4}$ on polyhedra and $n^{-1/2}$
on strongly convex level sets, that the choice between SPO$^+$ and PG is a
choice between a stronger guarantee under realisability and a weaker guarantee
that survives misspecification, and that the gain over the predict-then-optimise
baseline is empirically large in the misspecified regime that economic
forecasting actually occupies.
